% =============================================================================
%  strategy_D_dead_end.tex
%
%  A post-mortem on STRATEGY D (q_ramification_along_VF.tex): "re-resolve so
%  that eta_{C'} becomes a codimension-one point for q and apply the
%  codimension-one lemma iteratively".
%
%  RESULT: Strategy D cannot work, and the reason is a one-line invariant
%  computation:
%
%        e( V_{C'} | V_F ) = p^{r-1}  > 1   (r >= 2),
%
%  a WILD ramification index which is a birational invariant of the fixed
%  divisorial place V_{C'}|V_F.  No resolution can remove it, so q is never
%  unramified there and the codimension-one lemma (which only ever sees e=1)
%  cannot reach it.
%
%  As a bonus the note answers the natural follow-up: a wild q-defect of index
%  p^{r-1} does NOT prevent absorption of E/F -- on the contrary, the wildness
%  is exactly what ENABLES it (it multiplies the break to p^{r-1}*beta_E, which
%  is divisible by p, opening Artin-Schreier reduction over K_{C'}).  This
%  reframes the real target for Strategies A/B/C.
%
%  Compile:  TEXINPUTS=../..: latexmk -pdf strategy_D_dead_end.tex   (biber)
% =============================================================================
\documentclass[11pt]{article}
\usepackage{geometry}\geometry{margin=1in}
\usepackage[T1]{fontenc}
\usepackage{ifthen}
\usepackage{booktabs}

\setlength{\parindent}{1em}
\setlength{\parskip}{0.5em}

\input{includes}
\input{MacrosAndFigures}

\providecommand{\PP}{\mathbb{P}}
\providecommand{\fm}{\mathfrak{m}}
\providecommand{\Bl}{\mathrm{Bl}}

\usepackage{csquotes}
\usepackage[
backend=biber,
style=alphabetic,
sorting=ynt
]{biblatex}
\addbibresource{Bib.bib}

\newboolean{ThesisIntro}
\setboolean{ThesisIntro}{false}

\title{Strategy D is a dead end\\[2pt]
       {\large the $q$-defect at $V_F$ has wild ramification index $p^{r-1}$,
        a birational invariant; and why wildness \emph{enables} absorption}}
\author{Assaf Marzan}
\date{\today}

\begin{document}
\maketitle

\noindent\textbf{Purpose.}
Strategy D of \texttt{q\_ramification\_along\_VF.tex} proposed to dodge the
collision point by choosing a different resolution
$X_{C'}'\to C'^{(d)}$ in which $\eta_{C'}$ is a codimension-one centre for
$q:(\cP_{G/H})^{(d)}\to(\cP_{G/H})^{[d]}$, and then to apply the
codimension-one unramifiedness lemma
\Cref{lemma:torsor_unramified_codim_one} iteratively. This note shows the idea
cannot succeed, by a single invariant computation, and then settles the
follow-up question of whether the (necessarily) wild $q$-defect can still absorb
the ramification of $E/F$.

\medskip
\noindent\textbf{Summary.}
\begin{itemize}
  \item \Cref{thm:index}: $e(V_{C'}\mid V_F)=p^{r-1}$. For $r\ge2$ this is
        wild ($p\mid e$), so $q$ is wildly ramified along $V_F$.
  \item \Cref{cor:Ddead}: a ramification index of a fixed divisorial place is a
        birational invariant; no resolution changes $p^{r-1}$ to $1$. Strategy D
        is impossible, and the codimension-one lemma --- which only ever
        produces $e=1$ --- can never reach $V_{C'}$. This upgrades observation
        (O2) of the parent note from heuristic to theorem.
  \item \Cref{sec:absorb}: the wild index $p^{r-1}$ does \emph{not} block
        absorption; it is the mechanism that \emph{permits} it. The real
        obstruction is not the index but a finer, leading-term coincidence,
        which is what the global hypotheses (B)+(C) must rule out.
\end{itemize}

% #############################################################################
\section{The setup and the model}\label{sec:setup}
% #############################################################################

Notation as in \texttt{a\_wrong\_path.tex} and the parent note: after the
standard reductions ($k=\bar k$, $\cP\to U$ totally ramified at $p_\infty$,
\Cref{theorem:ramification_after_blowup_is_tame_for_p_power}), $G=\Z/p^r$,
$H\cong\Z/p$ the top layer, $G/H\cong\Z/p^{r-1}$ the bottom layer, $f:C'\to C$
the $G/H$-cover with $p'_\infty\mapsto p_\infty$. Write $\varpi$, $u$ for
uniformisers at $p_\infty$, $p'_\infty$; total ramification of the bottom layer
gives
\begin{equation}\label{eq:totram}
\varpi \;=\; u^{\,p^{r-1}}\cdot(\text{unit}) .
\end{equation}

We use the \emph{working} local model of the proven $\Z/p$ case
(\texttt{RamificationAfterBlowupIsTame.tex}, eq.~\cref{eq:complete_field_at_generic_point}
and the lines around it), \emph{not} the local models of the post-mortems
\texttt{a\_wrong\_path.tex} / \texttt{inertia\_symmetry\_obstruction.tex} (those
record dead ends and use mismatched valuations; e.g.\ the ``$J\supseteq G^d$''
claim there would already contradict the proven $\Z/p$ base case). In the
working model the exceptional-divisor valuation of the blowup of a $d$-fold
symmetric power at the collision point has an \emph{explicit uniformiser}: the
top elementary symmetric function. Thus
\begin{equation}\label{eq:unifs}
V_C:\quad \text{uniformiser } e_d=\textstyle\prod_{i=1}^d \varpi_i,\quad v_C(e_d)=1;
\qquad
V_{C'}:\quad \text{uniformiser } \epsilon_d=\textstyle\prod_{j=1}^d u_j,\quad v_{C'}(\epsilon_d)=1,
\end{equation}
where $\varpi_i$ (resp.\ $u_j$) are the base (resp.\ branch) coordinates of the
$d$ points. (Legitimate by the isomorphic-completion principle of
\texttt{RamificationAfterBlowupIsTame.tex}: only
$\widehat{\Oo}_{C',p'_\infty}\cong k[[u]]$ enters.)

% #############################################################################
\section{The ramification index of the \texorpdfstring{$q$}{q}-defect}\label{sec:index}
% #############################################################################

\begin{theorem}\label{thm:index}
$e(V_{C'}\mid V_F)=e(V_{C'}\mid V_C)=p^{r-1}$. In particular, for $r\ge2$ the
extension $K_{C'}/F$ is \emph{wildly} ramified at $V_{C'}\mid V_F$.
\end{theorem}

\begin{proof}
By \cref{eq:totram}, $\varpi_i=u_i^{\,p^{r-1}}\cdot(\text{unit})$ for each of the
$d$ branches, so the base uniformiser of $V_C$ pulls back to
\[
e_d=\prod_{i=1}^d\varpi_i
   =\Big(\prod_{j=1}^d u_j\Big)^{\!p^{r-1}}\!\!\cdot(\text{unit})
   =\epsilon_d^{\,p^{r-1}}\cdot(\text{unit}).
\]
Hence $v_{C'}(e_d)=p^{r-1}\,v_{C'}(\epsilon_d)=p^{r-1}$ by \cref{eq:unifs}.
Since $e_d$ is a uniformiser of $V_C$ and $\eta_{C'}\mapsto\eta_C$ under
$f^{(d)}$ (so $V_{C'}\mid V_C$), the ramification index is
\[
e(V_{C'}\mid V_C)=v_{C'}\big(\text{uniformiser of }V_C\big)=v_{C'}(e_d)=p^{r-1}.
\]
Finally $F/K_C$ is unramified at $V_C$ --- this is hypothesis (H1), equivalently
the $\Z/p^{r-1}$-symmetrised datum $\sum_i b(\varpi_i)$ lies in the residue field
$\kappa_C$ by the bound $m<d$ (the lemma of the working $\Z/p$ proof) --- so
$e(V_F\mid V_C)=1$ and
$e(V_{C'}\mid V_F)=e(V_{C'}\mid V_C)/e(V_F\mid V_C)=p^{r-1}$.
\end{proof}

\begin{rem*}
The same computation explains why \Cref{lemma:torsor_unramified_codim_one}
succeeds at boundary \emph{divisors} and only there. At the generic point of a
boundary divisor $Z_x$ exactly \emph{one} of the $d$ points sits at
$p'_\infty$ and the other $d-1$ are free; only one factor $u_j$ contributes to
the collision, the symmetric structure retains $d-1$ free copies of $K(C')$ in
the residue field, and one finds $e=1$ (trivial inertia). At $V_{C'}$
\emph{all $d$} points collide through the degree-$p^{r-1}$ wild cover, and
\cref{eq:totram} forces $e=p^{r-1}$. The codimension-one lemma is therefore not
``too weak to reach $V_{C'}$''; its conclusion is simply \emph{false} there.
\end{rem*}

% #############################################################################
\section{Why Strategy D cannot be repaired}\label{sec:Ddead}
% #############################################################################

\begin{cor}\label{cor:Ddead}
There is no birational model on which $q$ is unramified at the place
$V_{C'}\mid V_F$, and no factorisation of $q$ into codimension-one steps to which
\Cref{lemma:torsor_unramified_codim_one} applies and whose composite recovers
$V_{C'}$. Strategy D is impossible.
\end{cor}

\begin{proof}
The ramification index $e(w\mid w_F)$ of a fixed place $w=V_{C'}$ of $K_{C'}$
over its restriction $w_F=V_F$ to $F$ is an invariant of the \emph{valuations}
$w,w_F$ alone: it equals $[\,|w(K_{C'}^\times)| : |w_F(F^\times)|\,]$, the index
of value groups, and does not refer to any model. (Equivalently: $w$ and $w_F$
are intrinsic points of the Riemann--Zariski spaces of $K_{C'}$ and $F$; a
resolution only re-presents $w$ as a divisor, it cannot alter $w$.) By
\Cref{thm:index} this index is $p^{r-1}>1$. Any application of
\Cref{lemma:torsor_unramified_codim_one} yields a place with $e=1$; a
composite of such steps landing on $V_{C'}$ would give $e(V_{C'}\mid V_F)=1$,
contradicting \Cref{thm:index}.
\end{proof}

\noindent So the only thing a resolution can buy is to \emph{realise} $V_{C'}$
as a nice divisor for \emph{computing} its higher ramification (a service to
Strategy A), never to make $q$ unramified there. Strategy D, as a route to the
theorem, is closed.

% #############################################################################
\section{Does a wild $q$-defect still absorb $E/F$? Yes --- that is the point}
\label{sec:absorb}
% #############################################################################

It is tempting to hope that a \emph{large} wild defect $e=p^{r-1}$ is ``too
violent'' to let (H2) coexist with a ramified $E/F$, i.e.\ that wildness blocks
absorption. The opposite is true: wildness is the \emph{enabling} mechanism.

\paragraph{The mechanism.}
Present the top layer locally as Artin--Schreier over $\widehat F$,
\[
\widehat E=\widehat F(y),\qquad y^p-y=a,\qquad v_F(a)=-\beta_E,\ \ \gcd(\beta_E,p)=1,
\]
so $E/F$ is (if ramified) totally wildly ramified of break $\beta_E\ge1$
(\Cref{theorem:artin_schreier_ramification}(3)). Base-change along
$\widehat F\hookrightarrow\widehat{K_{C'}}$, of ramification index
$e=p^{r-1}$ (\Cref{thm:index}). The \emph{same} element $a$ now has
\begin{equation}\label{eq:scale}
v_{C'}(a)=e\cdot v_F(a)=-\,p^{r-1}\beta_E,
\end{equation}
and the pole order $p^{r-1}\beta_E$ is \textbf{divisible by $p$} (as $r\ge2$).
Divisibility by $p$ is exactly the trigger for Artin--Schreier reduction over
$\widehat{K_{C'}}$: writing $\pi$ for a uniformiser of $\widehat{K_{C'}}$ and
$a=c\,\pi^{-p^{r-1}\beta_E}+\cdots$ with $c=c_0^{\,p}$ ($k$ algebraically
closed), the substitution $y\mapsto y-c_0\pi^{-p^{r-2}\beta_E}$ replaces $a$ by
\[
a-\wp\!\big(c_0\pi^{-p^{r-2}\beta_E}\big)
 \;=\; a-c_0^{\,p}\pi^{-p^{r-1}\beta_E}+c_0\pi^{-p^{r-2}\beta_E},
\]
killing the leading pole. Iterating, the pole order can descend
$p^{r-1}\beta_E\to p^{r-2}\beta_E\to\cdots$, and \emph{if the lower-order data
cooperate it can reach $\le0$}, i.e.\ $\widehat E\widehat{K_{C'}}/\widehat{K_{C'}}$
becomes unramified --- (H2) holds while $E/F$ is ramified. This is precisely the
coincidence $\widehat{K_{C'}}=\widehat E$ engineered in
\Cref{ex:p2-counterexample}. Over $\widehat F$ no such reduction is available,
because there $v_F(a)=-\beta_E$ is \emph{prime to $p$} and $\widehat F$ has no
elements of the intermediate $\pi$-valuations $-p^{j}\beta_E$ ($j<r-1$) used by
the substitutions; the wild base change is exactly what supplies them.

\begin{quote}\itshape
The index $p^{r-1}$ neither causes nor prevents absorption by itself; it is the
gate that makes absorption arithmetically possible. Absorption \emph{completes}
only when the reduction terminates at a non-positive pole, which is a
\emph{leading-coefficient} condition, equivalent to $\widehat E\subseteq
\widehat{K_{C'}}$.
\end{quote}

\paragraph{Consequences for the live strategies.}
\begin{enumerate}
  \item The correct target is \emph{not} ``$\beta_q<\beta_E$'' as a bare
        inequality of indices, but the finer statement that the reduction of $a$
        over $\widehat{K_{C'}}$ \textbf{halts with a positive break}, equivalently
        $\widehat E\not\subseteq\widehat{K_{C'}}$. Generically the reduction
        halts at break $\beta_E$ (the prime-to-$p$ tail $c_0\pi^{-\beta_E}$
        cannot be removed), so $E\cdot K_{C'}/K_{C'}$ \emph{stays ramified} and
        (H2) forces the Goal; absorption is the non-generic, conspiratorial
        case.
  \item This is exactly Strategy C of the parent note, now with a precise
        criterion: show the global hypotheses
        \[
        \textbf{(B)}\ \mathrm{ram}(\cP)\le d,\qquad \textbf{(C)}\ r<d
        \]
        force the top-layer Artin--Schreier class $[a]\in
        \widehat F/\wp(\widehat F)$ to remain nontrivial after the base-change
        $\widehat F\hookrightarrow\widehat{K_{C'}}$, i.e.\ to survive in
        $\widehat{K_{C'}}/\wp(\widehat{K_{C'}})$. The class $[a]$ is built from
        the \emph{top} Witt component $a_{r-1}$ (break $\beta_E\sim u_r$), while
        $\widehat{K_{C'}}$ is built from the \emph{bottom} cover (breaks
        $u_1,\dots,u_{r-1}$); by Hasse--Arf $u_{r-1}<u_r$, and (B) keeps all
        $u_i<d$ while (C) keeps the number of layers below the blow-up scale.
        The task is to convert this separation of sources into non-containment
        of the local extensions.
  \item Strategy A is unaffected and remains the way to \emph{produce} the
        number: compute the reduced break of $a$ over $\widehat{K_{C'}}$ directly
        in the symmetric-power model. \Cref{thm:index} gives its starting input
        ($v_{C'}(a)=-p^{r-1}\beta_E$) for free.
\end{enumerate}

% #############################################################################
\section{Takeaway}\label{sec:takeaway}
% #############################################################################
\begin{itemize}
  \item \textbf{Strategy D is dead}, by the invariant $e(V_{C'}\mid V_F)=p^{r-1}$
        (\Cref{thm:index,cor:Ddead}). A fixed place cannot be de-ramified by
        re-resolving.
  \item \textbf{Wildness enables, not blocks, absorption} (\Cref{sec:absorb}):
        the base change multiplies the break to $p^{r-1}\beta_E\equiv0\!\!\pmod p$,
        opening Artin--Schreier reduction. The genuine obstruction lives one
        order deeper, in the leading coefficients --- i.e.\ in whether
        $\widehat E\subseteq\widehat{K_{C'}}$.
  \item \textbf{The fight is over a class in $\widehat F/\wp(\widehat F)$}, not
        over a ramification index. (B)+(C) must keep the top-layer class
        nontrivial over $\widehat{K_{C'}}$; the small decisive computation of the
        parent note ($r=2$, $d=3$) is precisely the test of this.
\end{itemize}

\printbibliography

\end{document}
